%% RESUMO EM LINGUA ESTRANGEIRA (elemento obrigatorio -- NBR 6028)
%% Traducao fiel do resumo em portugues. As keywords sao definidas em
%% config/dados.tex.

\begin{abstract}
This dissertation proposes and applies a multidimensional model to evaluate
document contextualization strategies in retrieval-augmented generation (RAG)
systems used to query Brazilian public institutional documents. The research
problem is that public agencies choose among strategies such as fixed-size,
recursive, semantic and structure-aware parent-child chunking without an
instrument that compares answer quality, cost and source traceability at once.
The model combines five dimensions --- retrieval, generation, efficiency,
traceability and institutional adequacy --- and its main contribution is the
Document Traceability Index (IRD), made of six binary criteria and validated
through agreement between independent raters and through convergent validity
with experts from public agencies. The model is applied in a paired experiment
with four strategies, a metadata ablation and a set of 60 to 100 reference
questions, keeping the generator model constant. As a technical product, the
research delivers a decision matrix for the agency's data and artificial
intelligence governance committee, evaluated for efficacy and formatively with
representatives of its user profiles.
\end{abstract}
